% Generated draft: requires mathematical and source review.
\documentclass[11pt]{article}
\usepackage[margin=25mm]{geometry}
\usepackage{fontspec}
\setmainfont{Latin Modern Roman}
\usepackage{amsmath,amssymb,mathtools,unicode-math}
\setmathfont{Latin Modern Math}
\usepackage{enumitem,needspace,xurl,hyperref}
\setlength{\parindent}{0pt}
\setlength{\parskip}{5pt}
\setlength{\emergencystretch}{3em}
\newcommand{\N}{\mathbb{N}}
\newcommand{\Z}{\mathbb{Z}}
\newcommand{\Q}{\mathbb{Q}}
\newcommand{\R}{\mathbb{R}}
\newcommand{\C}{\mathbb{C}}
\newcommand{\F}{\mathbb{F}}
\newcommand{\A}{\mathbb{A}}
\newcommand{\PP}{\mathbb{P}}
\newcommand{\E}{\mathbb{E}}
\newcommand{\eps}{\varepsilon}
\newcommand{\dd}{\,\mathrm{d}}
\begin{document}
\section*{UnsolvedMath ID: 700445}
\textbf{Source problem code:} catalogue entry
\subsection{Definitions and mathematical statement}
Let \(\Omega\) be a nonempty set and let \(G\leq\operatorname{Sym}(\Omega)\) be a permutation group. Its action is transitive if, for all \(\alpha,\beta\in\Omega\), some \(g\in G\) satisfies \(g\alpha=\beta\). A transitive action is primitive if its only \(G\)-invariant equivalence relations on \(\Omega\) are equality and the universal relation. The source does not specify whether degree-one actions are included in the term primitive; this boundary convention remains unresolved.

For \(\alpha\in\Omega\), define the point stabilizer \(G_\alpha=\{g\in G:g\alpha=\alpha\}\). A suborbit is an orbit \(\Delta=G_\alpha\beta\) of \(G_\alpha\) on \(\Omega\), where \(\beta\in\Omega\). The action of \(G_\alpha\) on \(\Delta\) is regular if for every \(\delta,\varepsilon\in\Delta\) there is exactly one \(h\in G_\alpha\) such that \(h\delta=\varepsilon\). This condition also makes that action faithful, meaning that only the identity fixes every point of \(\Delta\).

\textbf{Problem (A. N. Fomin, Kourovka Notebook 7.51).} Determine all primitive permutation groups, finite and infinite, for which there exist \(\alpha,\beta\in\Omega\) such that \(G_\alpha\) acts faithfully and regularly on \(G_\alpha\beta\). Equivalently, determine those admitting \(\alpha,\beta\in\Omega\) with
\[
G_\alpha\cap G_\beta=\{1\}.
\]
No finiteness assumption is imposed on the group or its permutation domain. The source imposes no requirement that \(\beta\neq\alpha\), and no such requirement is added here.
\subsection{Short English statement}
Determine which finite and infinite primitive permutation groups have a point stabilizer acting faithfully and regularly on at least one of its orbits. Regularity means that exactly one stabilizer element carries any chosen point of that orbit to any other. The convention concerning degree-one primitive actions requires clarification.
\subsection{Sources}
\begin{itemize}
\item Supplied local UnsolvedMath registry record: permanent ID \(700445\), exact source code \texttt{catalogue entry}, title \emph{Kourovka Notebook Problem 7.51}. The registry labels the problem open; that label is supplied metadata, not an independently verified assessment.
\item \emph{The Kourovka Notebook (Unsolved Problems in Group Theory)}, Problem 7.51, proposed by A. N. Fomin, seventh issue (1980). The supplied background identifies the reference edition as the twenty-first issue edition (2026), \texttt{arXiv:1401.0300v48}. The supplied collection description names E. I. Khukhro and V. D. Mazurov as editors. These bibliographic details are reproduced from the record without independent source verification.
\end{itemize}
\end{document}
